% Small-budget regime (C2 sweep), drafted 2026-09-05 from the plain-prompt sweep; chat-template numbers added when feat-005 closed. Prose revised 2026-09-05.
\subsection{The small-budget regime}\label{sec:sweep}

The released runs cover $k \in \{1, 3, 5\}$. To observe the mechanism where it acts, we reran the same six prompt classes (900 prompts, three trajectories each, $T_{\max}=200$, temperature 1) at $k \in \{0.1, 0.15, 0.25, 0.5, 1\}$, with the risky model alone and the anchor alone ($k=0$) as baselines on the same prompts and seeds. Each run was made twice, once with the raw prompt text and once with the risky model's chat template, the conversation markup it was instruction-tuned with (Figure~\ref{fig:sweep}). With raw prompts the constraint is active on $68$--$74\%$ of decode steps for $k \le 0.25$, on $49\%$ at $k=0.5$, and on $5.5\%$ at $k=1$; the mechanism becomes the identity only above $k=1$. No generation at $k \le 0.25$ is identical to the risky model's on the same seed, and one of 2{,}700 is at $k=0.5$, against $4\%$ at $k=1$ and $57\%$ between $k=3$ and $k=5$.

The anchor writes the opening of every answer. The mean number of leading tokens forced to the anchor is $39$ at $k=0.1$ (the released implementation's default is $k=0.15$, where it is $27$), $16$ at $k=0.25$, $8$ at $k=0.5$, and $4$ at $k=1$, matching $\lfloor \delta_{\mathrm{init}}/k \rfloor$. At the budgets where the certificate of Section~\ref{sec:certificate} is informative, a fifth of a $200$-token answer to a CopyBench prompt is written by a 1.8B model that has not read the book, and the fraction of steps forced to the anchor over the whole generation is $27\%$ at $k=0.1$. The copying metrics against the reference passages are flat across every budget and both baselines (ROUGE-L $0.058$--$0.062$, near-verbatim recall $0$), because the audited risky model does not reproduce the passages; a sweep of $k$ on this workload measures utility loss rather than protection, which is why Section~\ref{sec:attack-results} uses a memorising model.

The event-level tails are heavier at small budgets. On attack-split prompts the realised log-likelihood ratio exceeds $K$ for $13\%$ of trajectories at $k=0.1$, $16\%$ at $k=0.15$, $23\%$ at $k=0.25$, and $33\%$ at $k=0.5$, before falling to $10\%$ at $k=1$, with 95\% anytime-valid intervals that exclude zero at every budget (Figure~\ref{fig:llrtails}). The budgets at which the certificate is informative are thus the budgets at which up to a third of the outputs individually exceed it.

With the risky model's chat template the picture is the same but shifted. The instruct model's answers are shorter (89 tokens on average against 200, because it stops at its end-of-turn token); the anchor, which has never seen the template, is forced on $44\%$ of steps at $k=0.1$ and $31\%$ at $k=0.15$; the constraint is still active on $17\%$ of steps at $k=1$; and the anchor writes the first $35$ tokens at $k=0.1$ and $5.5$ at $k=1$. The realised log-likelihood ratio exceeds $K$ for $3$--$19\%$ of attack-split trajectories across the sweep. There were no invariant violations in any of the $37{,}800$ trajectories of the sweep.
